\section{Form degree in higher interaction terms}

\begin{proposition}[Quartic expressions on a three-manifold]
\label{prop:cbc26-quartic-degree}
Let $M$ be a smooth three-manifold and $A\in\Omega^1(M;W)$ for a finite-dimensional vector space $W$.
Any expression obtained by applying a four-linear map to coefficients of $A$ and wedging all four one-form factors is zero.
If a vector field $v$ and a four-linear map $T:W^{\otimes4}\to\mathbb C$ are specified, the expression
\[
 T(\iota_vA,A,A,A)\in\Omega^3(M)
\]
is instead a defined three-form.
It can be nonzero, but depends on the additional contraction and tensor.
\end{proposition}
\begin{proof}
The first expression belongs to $\Omega^4(M)=0$.
For the second, $\iota_v A$ has form degree zero, so the total degree is three.
For an explicit local example take $W=\mathbb C^3$, $A=e_1dx+e_2dy+e_3dz$, $v=\partial_x$, and the tensor whose only nonzero coordinate entry is $T(e_1,e_1,e_2,e_3)=1$.
Then the expression is $dx\wedge dy\wedge dz$.
\end{proof}

An RTT relation is an algebraic relation on matrix-valued series.
To use it as an action term, a map to a three-form such as the displayed contraction must be defined and its required symmetries proved.
A critical affine level alone does not make the Lie bracket in $\Omega^\bullet(M)\otimes\mathfrak g$ vanish.
For example, $A=e\,dx+f\,dy$ in $\mathfrak{sl}_2$ has $[A,A]/2=h\,dx\wedge dy\ne0$, independently of any level parameter.

\begin{proposition}[Integration on a finite quotient]
\label{prop:cbc26-quotient-integral}
Let a finite group $\Gamma$ act freely by orientation-preserving diffeomorphisms on a closed oriented three-manifold $M$.
Let $\omega$ be an invariant three-form and $\bar\omega$ its descended form on $M/\Gamma$.
Then
\[
 \int_{M/\Gamma}\bar\omega=\frac1{|\Gamma|}\int_M\omega.
\]
\end{proposition}
\begin{proof}
The quotient map is an oriented covering of degree $|\Gamma|$.
An evenly covered coordinate cover and a partition of unity reduce integration to the sum over its $|\Gamma|$ sheets.
Each sheet has the same integral, giving the formula.
\end{proof}

For a nonfree finite action, integration on the quotient stack uses the right-hand side as its normalization.
An additional division by $|\Gamma|$ after taking this quotient integral changes the normalization.
Neither a fractional grading nor a W-algebra specifies an action of $\Gamma$ on $M$.
